
\chapter{Fallbeispiele}

\ChapterEnd

\chapter{Arztbriefe}

\section{Chirurgische Arztbriefe}

\subsubsection{Gallenkolik}


Wir berichten über unsere Patientin  Frau Anna Musterfrau,
geboren am 30.03.1969, welche in der Zeit vom
01.05.2013 bis 05.05.2013 an unserer Abteilung stationär
behandelt wurde.

\paragraph{Diagnose}
\begin{itemize}
  \item Choledocholithiasis\\
Abdominalgie
\end{itemize}


\paragraph{Nebendiagnosen}

\begin{compactitem}
  \item Chronisches Zervicalsyndrom
  \item Adipositas 
  \item Nikotinabusus 
  \item Arterielle Hypertonie 
  \item Zustand
nach Appendektomie 1993
\end{compactitem}

\paragraph{Therapie} Konservativ, Schmerz- und Infusionstherapie


\paragraph{Epikrise}

Die Patientin gelangt in den 
Nachtstunden über unsere Notfallaufnahme zur stationären
Aufnahme. Sie klagte über kolikartige Bauchschmerzen
mit punctum maximum im rechten Oberbauch seit den
frühen Abendstunden. Sie wurde noch in der Notaufnahme
anlagetisch mit Buscopan und Novalgin behandelt, es stellte sich eine deutliche
Beschwerdebesserung ein. Im durchgeführten Abdomen-Schall zeigte
sich eine deutliche Gallenabflussstörung sowie eine
Sludge-Gallenblase ohne Zeichen einer akuten Entzündung. 
Im Labor zeigten sich war keine
Entzündungszeichen, jedoch geringfügig erhöhte
Leberfunktionsparameter und Cholestasewerte. 

Mit der
Patientin wurde ausführlich das weitere Procedere
besprochen, in Anbetracht dessen das in  \Range{4}{5} Monaten eine
Fernreise geplant ist, entschließt sich die Patientin auch
für eine elektive Cholezystekotime. Es wird in OP-Termin
für den 7. Juni 2013 vereinbart. Es erfolgt noch während des Aufenthaltes
eine erste Aufklärung über den Eingriff und damit
verbundenen Risiken. Die Patientin kann sodann am 05.05.2013
beschwerdefrei aus der stationären Behandlung entlassen werden.

\paragraph{Weiteres Procedere} 

\begin{itemize}
  \item Kontrolle der Leberfunktionsparameter und
der Entzündungswerte einer Woche durch den Hausarzt
\item Auf Gellen-Schonkost und fettarme Ernährung ist zu
achten
\item 
Die erforderlichen Voruntersuchungen wurden Patienten
auf einem separaten Merkblatt mitgegeben. Diese
Untersuchungen bitte war bis zwei Wochen zwotem
Aufnahmetermin durchführen lassen.
\item Bitte Wiedervorstellung
eine Woche vor dem Aufnahmetermins in unserer chirurgischen
Ambulanz, vormittags acht bis 11:00 Uhr zur Erledigung der
Aufnahmeformalitäten und Vorstellung an der
Anästhesiologischen Abteilung.
\item Wiederaufnahme am
06.07.2013 auf der Station 3B um 8:30 Uhr mit allen
vorhandenen Befunden.
\end{itemize}


\paragraph{Medikation} 

Concor 5 mg Tabletten 1-0-1 \\
Amlodipin 10 mg
1-0-0\\
Buscapine plus Tbl. bei Bedarf bis zu dreimal täglich \\
Novalgin 500 mg Tbl. bei Bedarf bis zu dreimal täglich

Wiedervorstellung bei Beschwerden hierorts jederzeit möglich

Gezeichnet Dr. Schrammel 
Stationsarzt 


\section{Gynäkologie}

\subsection{Gynäkologische Patienten mit Hirnmetastasierung}

Briefkopf einfügen

Löblicher Kollege Rufzeichen

Wir berichten über die Patientin Ulrike Meyer geboren am
21.07.1959, welche in der Zeit vom 27.06 bis 18.07.2013 an
unserer Abteilung stationär aufgenommen war. 


\paragraph{Diagnose:}

\begin{itemize}
  \item Symptomatischer generalisierter epileptischer Krampfanfall
bei Verdacht auf ihren mittels Dosierung.
\end{itemize}



\paragraph{Nebendiagnosen:}
\begin{compactitem}
\item 
Histologisch nicht verifizierte Raumforderung im kleinen
Becken vermutlich von der Adnex-Region ausgehend 
\item Raumforderungen in der Leber, einer Metastase
entsprechend
  \item Adipositas 
  \item   depressive Episode 
  \item   Arterielle
Hypertonie 
\item Zustand nach Beinvenenthrombose rechts 1991
\item Zustand nach Myocardinfarkt 01/2001 
\item Koronare
Zweigefäßerkrankung (LAD,
CX) 
\item Zustand nach primärer PTCA + BMS LAD 01/2001
\item Zustand nach sekundärer PTCA + DES CX 05/2001
\end{compactitem}


\paragraph{Medikation} 
Concor
5 mg 1-0-1 \\
Amlodipin 10 mg 1-0-0 \\
Halcion Tbl. 0,25 mg 
0-0-0-\nicefrac{1}{2}\\ 
Lasix 40 mg Tabletten 1-0-0
\\
ThromboASS 100 mg pausiert 
\\
Novalgin Tropfen 25 gtt bis zu
dreimal täglich bei Bedarf 
\\
Keppra 500 mg 1-1-1 Revision Anm.
bitte Dosierung nach sehen

\paragraph{Kurzdekurs} Die Patientin kommt ursprünglich zur geplanten
Aufnahme zum Staging einer histologisch nicht
verifizierten Raumforderung. Während des Aufenthaltes
erlitt die Patientin einen tonisch-klonischen Krampfanfall und
wurde umgehend auf unsere Intensivstation verlegt, dort
stabilisierte sich ihr Zustand umgehend. Es wurde
eine Computertomographie des Gehirns durchgeführt, bei
dieser zeigte sich eine bisher noch nicht bekannte
Raumforderungen desTemporallappen rechts. Der
beigezogenen Konsiliar-Neurologe konnte keine bleibenden
neurologischen Defizite feststellen, sie wurde mit
Keppra und ??? zu Anfallsprophylaxe eingestellt. Die Patientin
konnte am nächsten Tag wieder die Normalstation verlegt
werden. Es wurden wiederholt war Gespräche mit der Patientin
zu Ihrer Situation geführt, ebenso wurde eine psychologische
Betreuung veranlasst. Mit der Patienten wurde offen über die
 möglichen Konsequenzen gesprochen,  sie
verweigert weiterhin den histologische Abklärung der
Raumforderung, ebenso wünscht sie ausdrücklich keine kausale Therapie im
Sinne mit Chemo- oder Strahlentherapie. Es wird somit mit der
Patientin eine Best Supportive Care vereinbart, eine
Vorstellung bei unserem Palliativ-Team ist erfolgt. Die
Patientin kann sodann im guten Allgemeinzustand aus der
stationären Behandlung entlassen werden. Ein Termin für eine
Wiedervorstellung wurde vereinbart. 

\paragraph{Weiteres Procedere} 
\begin{itemize}
  \item Wiedervorstellung bei Beschwerden hierorts jederzeit möglich.
  \item 
Vorstellung beim niedergelassenen Neurologen sobald möglich.
Kontrolle des Blutbilds und Leber.Funktionswerte in
einer Woche beim Hausarzt.
\item Wiedervorstellung zwei Wochen in unserer gynäkologischen
ambulanzvormittags erbeten.
\end{itemize}



Dr. Schrammel\\
Sekundararzt


\subsection{Geburt}


Briefkopf einfügen

Martina Musterfrau geboren 30.3.1982

\paragraph{Diagnose} Gravida II Partus II SSW 38+6 

\paragraph{Therapie}
Spontane vaginale Geburt eines reifen Einlings (männlich,
Florian, 3120 g, am 05.05.2013 um 00:05 Uhr)

\paragraph{Therapie} 
Aktiferrin Kps. 1-0-0

\paragraph{Procedere} 

\begin{enumerate}
  \item Gynäkologisch-fachärztliche Kontrolle
in \Range{6}{8} Wochen 
\item 
Vorstellung in unserer Kinderambulanz
in zwei Tagen 
\item 
Vorstellung der niedergelassenen Orthopäden in vier Wochen
\end{enumerate}



Dr. Meyer\\
Sekundararzt 




\ChapterEnd